% One bit, divided into time quanta, in both phases, with the arithmetic.
\begin{tikzpicture}[font=\sffamily\scriptsize, x=1cm, y=1cm]

\tikzset{
  tq/.style={draw=linecol, anchor=south west, minimum height=6mm,
             minimum width=6mm, inner sep=0pt, font=\sffamily\tiny},
  tqsync/.style={tq, fill=hwcol},
  tqseg1/.style={tq, fill=kercol},
  tqseg2/.style={tq, fill=usrcol},
}

% ================================================================ the nominal bit
\node[signame] at (-0.3,2.9) {one nominal bit};
\node[tqsync] at (0,2.6) {S};
\foreach \i in {0,...,9}{ \node[tqseg1] at ({0.6+\i*0.6},2.6) {}; }
\foreach \i in {0,...,2}{ \node[tqseg2] at ({6.6+\i*0.6},2.6) {}; }
\draw[deadline] (6.6,2.4) -- (6.6,3.6);
\node[tlabel, anchor=south, text=red!60!black] at (6.6,3.6) {the sample point};
\node[tlabel, anchor=north] at (0.3,2.55) {sync};
\node[tlabel, anchor=north] at (3.6,2.55) {segment 1: propagation and phase};
\node[tlabel, anchor=north] at (7.5,2.55) {segment 2};
\node[chlabel, anchor=west] at (9.0,2.9) {14 time quanta, sample at 11 of 14: 79 per cent};

% ================================================================ the data bit
\node[signame] at (-0.3,1.1) {one data bit};
\node[tqsync, minimum width=3mm] at (0,0.8) {};
\foreach \i in {0,...,8}{ \node[tqseg1, minimum width=3mm] at ({0.3+\i*0.3},0.8) {}; }
\foreach \i in {0,...,2}{ \node[tqseg2, minimum width=3mm] at ({3.0+\i*0.3},0.8) {}; }
\draw[deadline] (3.0,0.6) -- (3.0,1.8);
\node[chlabel, anchor=west] at (9.0,1.1) {13 quanta at a shorter quantum: 2 Mbit/s};

\node[chlabel, anchor=west] at (4.2,0.95) {the same structure, a quarter of the time};

% ================================================================ the arithmetic
\node[chlabel, anchor=north west] at (-0.3,0.2) {the arithmetic, written out because the calculator that does it is under a network copyleft licence};

\node[regcell, minimum width=104mm, text width=102mm, anchor=north west,
      fill=white, align=left, minimum height=20mm] at (-0.3,-0.3)
  {\ttfamily\scriptsize
   quantum = kernel\_clock / prescaler\\
   bit     = 1 sync quantum + segment1 + segment2\\
   rate    = quantum / quanta\_per\_bit \hfill both divisions must be exact\\
   sample point = (1 + segment1) / quanta\_per\_bit \hfill 75 to 87.5 per cent\\
   resynchronisation jump width $\le$ segment2 \hfill and $\le$ the register's maximum};

% ================================================================ the two gotchas
\node[umlnote, text width=52mm, anchor=north west] at (-0.3,-3.0)
  {Each phase has its own prescaler and its own segments, so the two need not
   have the same number of quanta. What they must share is the kernel clock,
   which is why step 1 reads it back rather than assuming it, and why a rate
   that does not divide exactly stops the node instead of being rounded.};

\node[umlnote, text width=52mm, anchor=north west] at (5.8,-3.0)
  {Above about one megabit in the data phase, the round trip through the
   transceiver is a significant fraction of a bit, and the controller must
   compensate for it. The symptom of omitting that is precise: frames that pass
   in loopback and fail on a real bus, with an error rate that rises with the
   data rate.};
\end{tikzpicture}
